% Gerado a partir de docs/tcc.md por scripts/sincronizar_tcc.py.
\begin{otherlanguage*}{english}
\begin{resumo}[Abstract]
% Gerado a partir de docs/tcc.md por scripts/sincronizar_tcc.py.
This work presents the design and initial development of a contextual evaluation system for digital public services. The research problem concerns linking citizens\textquotesingle{} declared evaluations to the navigation context in which feedback is requested. The work combined bibliographic and regulatory review, Lean Inception, and incremental prototyping, drawing on Design Science Research. A local Python prototype was implemented with an API, simulated pages, operational rules based on clicks and reported duration, a feedback form, and an administrative dashboard. Telemetry provides behavioral indications and does not measure perceived quality. Stars and attributes capture declared judgments; this version retains NPS as an auxiliary recommendation question whose interpretation in public services has not been validated. Separately from the functional implementation, eight participants evaluated interface flows in Figma, with automatic click and time records collected through Useberry and a post-session questionnaire. Four participants performed each flow, with artificial-intelligence-assisted allocation described as a draw without an auditable record, and with a redirect popup present in both flows. The results describe task completion, perceived ease, and post-session reports. They do not establish greater participation, more specific feedback-form comments, correct API-trigger behavior, or superiority of contextual feedback. The investigative contribution is the comparative hypothesis and the distinction between behavioral and declared evidence; the technological contribution is the implementation that can connect them. Further work requires operational measures, a documented protocol, and rule calibration before deployment.

\vspace{\onelineskip}
\noindent \textbf{Keywords:} digital public services. contextual feedback. service evaluation. user experience. Design Science Research.
\end{resumo}
\end{otherlanguage*}
